The category of core methods as a layered Hasse diagram (\bgref{A.16}), from Frozen (initial, bottom) to Full FT (terminal, top); the dashed outer frame stands for the A1 arrows Frozen${}\to M\to{}$Full FT of every weight-space method $M$ (\thmref{Observation}{I.4}(a)). Chips are the \nCore{} core methods, coloured by modification kind and dashed for the $57$ that are not objects of the weight-space slice. An arrow $M\to N$ points up and means that $N$ simulates $M$ through a pointed smooth map $h$ with $\rho_N\circ h=\rho_M$ (\thmref{Definition}{I.2}). All $161$ recorded arrows between two core methods are drawn, coloured by rule (Table~\ref{tab:atlas-rules}): $29$ by freezing (A2, hooked tail), $122$ by factorisation (A3), $2$ by products (A4) and $8$ by sums or stages (A5). The $7$ loops land on another member of a family, as $\mathrm{LoRA}_r\to\mathrm{LoRA}_{r+1}$ does, and the $5$ classes of mutually simulating families share a rung in a grey box. Arrows into a method with four or more incoming arrows from outside its class end in one trunk ($57$ into LoRA), and the $28$ methods whose arrows all land in the LoRA class form the comb beneath it. The $26$ lighter arrows are composites, also reached by a longer path of arrows.
